\section{The chronological operator of a dyadic protocol}
\label{sec:peps_dyadic_protocol_circuit}

The changes of a supplied protocol determine its actual final vector before
any distributed compression is applied. The following identities record this
composition, including all private registers. They are the ordered products
used in \cite[Section~8.1]{OpenAI2026PolynomialPEPS}.

\begin{definition}[Protocol operator]
    \label{def:peps_dyadic_protocol_operator}
    \leanok
    \uses{def:peps_dyadic_protocol, def:peps_original_chronology_composite}
    \lean{TNLean.PEPS.Approximation.DyadicProtocol.originalCircuit}
    Let $W_t:\mathcal H_{A_t}\to\mathcal H_{A_{t+1}}$, $0\le t<T$,
    be the changes of a dyadic protocol. Let $I_0$ identify the initial memory
    with the product of the original party memories, and let $I_T$ identify
    the final memory with the physical registers followed by the private
    registers. The protocol operator is $I_T W^{[T]} I_0^{-1}$.
\end{definition}

\begin{theorem}[Evaluation on the actual trajectory]
    \label{thm:peps_dyadic_protocol_operator_eval}
    \leanok
    \uses{def:peps_dyadic_protocol_operator}
    \lean{TNLean.PEPS.Approximation.DyadicProtocol.eval_chain_trajectory}
    \lean{TNLean.PEPS.Approximation.DyadicProtocol.eval_originalCircuit}
    If $x$ is the original product input and
    $\Psi_0=I_0^{-1}x$, $\Psi_{t+1}=W_t\Psi_t$, then
    $W^{[t]}I_0^{-1}x=\Psi_t$ at every stage. In particular,
    $(I_T W^{[T]}I_0^{-1})x=I_T\Psi_T$.
\end{theorem}
\begin{proof}\leanok
    \uses{thm:peps_original_chronology_composition}
    The stage identity follows by induction from
    $W^{[0]}=1$ and $W^{[t+1]}=W_tW^{[t]}$.
    Applying $I_T$ to the final-stage identity gives the asserted output.
\end{proof}

\begin{theorem}[Physical density in site coordinates]
    \label{thm:peps_dyadic_protocol_density_coordinates}
    \leanok
    \uses{def:peps_dyadic_protocol, def:peps_dyadic_protocol_operator,
        def:peps_original_recorded_gate_maps}
    \lean{TNLean.PEPS.Approximation.DyadicProtocol.reindex_original_physicalDensity}
    \lean{TNLean.PEPS.Approximation.DyadicProtocol.rectangularTraceNorm_reindex_sub_physicalDensity}
    Let $\rho_{\mathrm{party}}$ be the physical density obtained by evaluating
    the original protocol circuit on its product input and tracing out its
    private registers. Let $E$ be the physical configuration bijection from
    party configurations to configurations of the square. Then the actual
    protocol density satisfies
    $\rho_{\mathrm{protocol}}(x,y)
        =\rho_{\mathrm{party}}(E^{-1}x,E^{-1}y)$ for every pair of square
    configurations $x,y$. Neither density is normalized in this identity.
    For any matrix $M$ indexed by party configurations, put
    $M^E(x,y)=M(E^{-1}x,E^{-1}y)$. Then
    \begin{align}
        \|M^E-\rho_{\mathrm{protocol}}\|_1
         & =\|M-\rho_{\mathrm{party}}\|_1.\notag
    \end{align}
\end{theorem}
\begin{proof}\leanok
    \uses{thm:peps_original_exact_gate_maps, thm:peps_dyadic_protocol_operator_eval}
    The recorded gate sums evaluate to the original circuit, whose output is
    the actual final trajectory. Write $u$ for that vector in party physical
    coordinates and private coordinates. In site coordinates its value at
    $(x,a)$ is $u(E^{-1}x,a)$. Thus
    \begin{align}
        \rho_{\mathrm{protocol}}(x,y)
         & =\sum_a u(E^{-1}x,a)\overline{u(E^{-1}y,a)} \notag \\
         & =\rho_{\mathrm{party}}(E^{-1}x,E^{-1}y).\notag
    \end{align}
    This is the partial-trace identity under the product of $E$ with the
    identity on the private coordinates. Relabelling the rows and columns
    by a bijection commutes with subtraction and preserves the trace norm,
    which gives the final identity.
\end{proof}
